\scanpage{144}
\uline{Construction abstraite d'une théorie de motifs semi-simples.}

On part d'une catégorie $\mathscr V$ \deleted{\unsure{pré}}``additive'', \deleted{\unsure{i.e. où}}\inserted{en particulier} les \deleted{Hom}\inserted{Fl}$(X,Y)$ \deleted{sont} munis d'une structure de groupe abélien, les accouplements de composition étant bilinéaires. \deleted{[On ne suppose pas qu'il existe un objet nul, ni des sommes \unsure{ou} produits de deux objets .....]} On s'intéresse, pour toute catégorie \uline{abélienne} $\mathcal C$, aux foncteurs \deleted{contravariants} ``additifs'' (i.e. commutant aux sommes finies, \unsure{ou ce qui revient au|ou encore} $=$, aux produits finis)
\[
 H\deleted{\unsure{F}}:\mathscr V^{\deleted{\unsure{\circ|*}}}\longrightarrow\mathcal C
\]
et on cherche un foncteur universel de cette nature
\[
 \underline{h}:\mathscr V^{\deleted{\circ}}\longrightarrow\mathcal M(\mathscr V)=\mathcal M
\]
foncteur additif dans catégorie abélienne $\mathcal M$) tel que pour toute catégorie abélienne $\mathcal C$, l'application
\[
\begin{aligned}
 T&\mapsto T\circ\underline h\\
 \operatorname{Fonc\,exacts}(\mathcal M,\mathcal C)&\longrightarrow\operatorname{Fonc\,add}(\mathscr V^{\deleted{\circ}},\mathcal C)
\end{aligned}
\]
soit bijective (ce qui caractériserait $\mathcal M$ à \deleted{\unsure{équivalence}} isom près), ou le cas échéant, que
\[
 \underline{\operatorname{Fonc\,exacts}}(\mathcal M,\mathcal C)\ \unsure{\mathord{-}|\longrightarrow}\ \underline{\operatorname{Fonc\,add}}(\mathscr V^{\deleted{\circ}},\mathcal C)
\]
soit une équivalence de catégories (ce qui
% archive sheet number: 143
